% Research fragment for integration by the parent agent.
% Requires amsmath, amssymb, amsthm and theorem/lemma/proof environments.
\section{Scalar iteration uniqueness already defines predecessor}
\label{sec:ur-predecessor}

This section gives the predecessor portion of the completion of the
scalar-uniqueness direction identified in the incoming
\texttt{scope\_kernel} report.  Its new ingredient is a modified scalar
ramp whose invariant is a direct fusion law.  Equality reflection is
proved independently in Section~\ref{sec:ur-reflection-completion}.

Work in ordinary intuitionistic first-order logic with actual equality.
Terms are generated by variables, $0$, $S$, and
$I(n,a;u.h)$, with only $u$ bound in $h$.  Assume the two iterator equations,
\[
 I(0,a;u.h)=a,\qquad I(Sn,a;u.h)=h[I(n,a;u.h)/u],
\]
and numerical coverage $n=0\lor\exists m\,n=Sm$.
Replace equation induction by the contextual scalar uniqueness rule
\[
 \frac{\Gamma\vdash f(Sz)=h[f(z)/u]}
 {\Gamma\vdash f(t)=I(t,f(0);u.h)}\quad(\mathrm{UR}),
\]
where $z$ is fresh for $\Gamma$ and the passive parameters, $f(z)$ is a
term template, and $h(u)$ has no free $z$.  All bodies are finite earlier
syntax.  A useful derived rule says that $f$ and $g$ are equal whenever
they have provably equal initial values and the same unconditional
autonomous update.  Apply UR to both and transport the initial value.
No induction hypothesis is permitted in the recurrence premise.

\begin{theorem}
The preceding theory defines a finite scalar term $P$ and proves
$P(0)=0$ and $P(Sn)=n$.  Consequently it proves successor injectivity.
The derivation does not use separation of $0$ and $S0$.
\end{theorem}

\subsection{Elementary identities by uniqueness}
Write $1=S0$, set $C(n,a,b)=I(n,a;u.b)$, $Z(n)=C(n,1,0)$, and
$a+b=I(b,a;u.Su)$.  Coverage gives $Z(n)=0\lor Z(n)=1$ and
\[
 C(0,a,b)=a,\quad C(Sn,a,b)=b,\quad
 C(e,Sa,Sb)=S(C(e,a,b)).
\]
The last identity follows directly by cases on $e$.
All commutative-monoid laws for addition follow from common-update
uniqueness.  In detail, $n$ has initial value $0$ and update $S$, so
$0+n=n$; $S(a+n)$ has initial value $Sa$ and update $S$, so
$Sa+n=S(a+n)$.  Thus $n+a$ has initial value $a$ and update $S$,
proving $n+a=a+n$.  Finally $a+(b+n)$ has initial value $a+b$ and
update $S$, proving associativity.  These facts justify the usual finite
sum notation $2a=a+a$ and $3a=a+a+a$.

For any fixed body $h$, the following are also common-update
uniqueness proofs:
\[
 I(m+n,a;h)=I(n,I(m,a;h);h),\qquad
 I(Sn,a;h)=I(n,h(a);h).
\]
For the first compare the two sides as functions of $n$, with update $h$
and initial value $I(m,a;h)$; for the second the initial value is $h(a)$.
A constant function fixed by $h$ is likewise the corresponding iterator.

\subsection{Two finite clocks}
Define
\[
 N(n)=I(n,0;u.Z(u)),\qquad j(u)=Z(u)+2N(u),\qquad
 M(n)=I(n,0;u.j(u)).
\]
Coverage on the counter of $N$ proves that $N(n)$ is $0$ or $1$.
Hence $Z(Z(N(n)))=N(n)$ and $N(n)+Z(N(n))=1$ by two finite cases.
The computation rule gives $N(n+2)=N(n)$.
The function $N(2n)$ has initial value $0$ and identity update, whereas
$N(3n)$ has initial value $0$ and update $Z$.  UR therefore gives
\[
 N(2n)=0,\qquad N(2n+1)=1,\qquad N(3n)=N(n).
\]

The range of $j$ lies in $\{0,1,2\}$: at $u=0$ its value is $1$;
at $u=Sv$ the term $Z(u)$ is zero and $2N(u)$ is $0$ or $2$.
Coverage on the counter of $M$ therefore proves that its range also lies
in $\{0,1,2\}$.  Direct computation gives
\[
 j(0)=1,\qquad j(1)=2,\qquad j(2)=0.
\]
Thus $M(n+3)=M(n)$ by three computations and the three cases for
$M(n)$.  UR with the identity update gives $M(3n)=0$, and one further
computation gives $M(3n+1)=1$.

\subsection{The modified ramp and its fusion law}
Use the update
\[
 g(u)=u+N(M(u))+1,\qquad G(n)=I(n,0;u.g(u)).
\]
It differs from the older ramp update $u+M(u)+1$ away from its intended
orbit.  That difference makes the following law valid at every argument:
\[
                         M(g(u))=Z(M(u)).                 \tag{1}
\]
Indeed, when $M(u)=0$, $g(u)=u+1$ and the new $M$ value is $1$;
when $M(u)=1$, $g(u)=u+2$ and the new value is $0$; when $M(u)=2$,
$g(u)=u+1$ and the new value is again $0$.
Consequently $M(G(n))$ has initial value $0$ and update $Z$.
Uniqueness, without any induction hypothesis, proves
\[
                             M(G(n))=N(n).                \tag{2}
\]
As $N$ fixes the two Boolean values, this gives
\[
 G(Sn)=G(n)+N(n)+1,\qquad G(n+2)=G(n)+3,
 \qquad G(n+3)=G(n)+N(n)+4.                               \tag{3}
\]
The last two equations follow by two or three computations and the
two cases for $N(n)$.

The functions $G(2n)$ and $3n$ both start at zero and have update
$v\mapsto v+3$, so UR yields $G(2n)=3n$.
For the less immediate scaling equation, put
\[
 f(n)=G(3n),\qquad k(n)=3G(n)+N(n),\qquad H_0(v)=v+M(v)+4.
\]
Equation (2) and parity scaling give $M(f(n))=N(n)$.
The identities $M(3a)=0$ and $M(3a+1)=1$, together with the two cases
for $N(n)$, give $M(k(n))=N(n)$.
Equations (3) show $f(Sn)=H_0(f(n))$.  In those same two Boolean cases,
\[
 \begin{split}
 k(Sn)&=3G(n)+3N(n)+3+Z(N(n))\\
      &=3G(n)+2N(n)+4=H_0(k(n)).
 \end{split}
\]
Both start at zero.  A second common-update uniqueness proof therefore
gives
\[
                  G(2n)=3n,\qquad G(3n)=3G(n)+N(n).       \tag{4}
\]

\subsection{The scalar comparison recurrence}
Set $H(u)=C(N(u),G(u),u)$.  From (4) and the parity identities,
\[
                 H(2u)=3u,\quad H(1)=1,\quad H(3u)=3H(u).
\]
For the last identity split $N(u)=0$ or $1$; in the zero case use
$G(3u)=3G(u)$, and in the one case both sides retain the corresponding
odd argument.  No division theorem is used.
The functions $I(y,3a;H)$ and $3I(y,a;H)$ have the same initial
value $3a$ and update $H$, hence
\[
 I(y,3a;H)=3I(y,a;H),\qquad I(y,1;H)=1.
\]
Define
\[
 E(x)=I(x,1;u.2u),\qquad L(x,y)=N(I(y,E(x);H)).
\]
The rotation identity and the displayed equations yield
\[
 L(0,y)=1,\qquad L(Sx,0)=0,\qquad L(Sx,Sy)=L(x,y).        \tag{5}
\]
For the final equality, rotate the first iterator step at the initializer
$2E(x)$, replace $H(2E(x))$ by $3E(x)$, move the factor three through
the iterator, and use $N(3v)=N(v)$.

\subsection{Saturated iteration and predecessor}
Define
\[
 B_m(u)=C(L(Su,m),u,Su),\qquad W(t,m)=I(t,0;u.B_m(u)).
\]
Equation (5) and selector compatibility with successor give
\[
 B_{Sm}(Su)=S B_m(u),\qquad B_{Sm}(0)=1.
\]
As functions of $t$, both $W(St,Sm)$ and $S W(t,m)$ start at one
and have update $B_{Sm}$.  Therefore UR proves
\[
                             W(St,Sm)=S W(t,m).           \tag{6}
\]
The function $n\mapsto L(Sn,n)$ starts at zero and has identity update
by (5), so it is identically zero by UR.  By the already universal
equation (6), $n\mapsto W(n,n)$ starts at zero and has update $S$,
so it equals $n$.  It follows that
\[
 W(Sn,n)=B_n(W(n,n))=C(L(Sn,n),n,Sn)=n.                   \tag{7}
\]
Finally put
\[
 P(m)=I\bigl(m,0;u.C(L(SSu,m),u,Su)\bigr).
\]
The zero equation is immediate.  For fixed $m$, the function
\[
 V(t,m)=I\bigl(t,0;u.C(L(SSu,Sm),u,Su)\bigr)
\]
has the same initial value and update as $W(t,m)$, because (5)
identifies its update pointwise with $B_m$.  UR yields $V(t,m)=W(t,m)$.
Using (7),
\[
 P(Sm)=V(Sm,m)=W(Sm,m)=m.
\]
Applying $P$ to $Sx=Sy$ now gives $x=y$.

\subsection{The exact dependency boundary}
Every appeal to induction in the predecessor bootstrap has been replaced
by a specified unconditional common-update uniqueness instance.  All case
distinctions use coverage and finite ranges; no separation axiom, equality
characteristic, or induction on an equation has been smuggled in.
One also obtains $d(x,y)=I(y,x;u.Pu)$ and
\[
 d(Sx,Sy)=d(x,y),\qquad d(a+n,n)=a,
\]
again by common updates $P$ and the identity, respectively.
Consequently addition is cancellative, and
$m(x,y,z)=d(x+z,y)$ is a Mal'cev operation:
$m(x,x,z)=z$ and $m(x,y,y)=x$.

The equality-reflection step needs a different argument.  The older
proof used
\[
 C\bigl(\operatorname{sg}(d(x,y)),x,y+d(x,y)\bigr)=x,
 \qquad \operatorname{sg}(v)=C(v,0,1).
\]
Its familiar proof invokes the equality at $y$ as an induction hypothesis.
The left side has a counter-dependent reset recurrence, not an immediate
autonomous scalar recurrence.  It is not imported into this bootstrap.
Instead, Section~\ref{sec:ur-reflection-completion} constructs an explicit
scalar accumulator using Boolean diagonal descent, proves a maximum
commutativity identity, and obtains actual equality reflection from it.
